% =====================================================================
% Appendix
% =====================================================================
\section{Implementation details}
\label{app:details}

\subsection{Compute and reproducibility}

Everything runs on CPU. The models are small, under $10^6$ parameters, and the
gauge relies on \texttt{matrix\_exp} and \texttt{linalg.lstsq}, neither of which
is implemented for Apple's MPS backend. Python, NumPy and PyTorch are seeded per
run and deterministic kernels are requested. The five-seed suite takes a few
hours on eight cores. A single entry point reproduces every number and figure
from scratch; see the README.

\begin{figure}[t]
  \centering
  \gndfig{fig2_architecture.pdf}{\textwidth}
  \caption{\textbf{Architecture and objective.} Only the gauge is context
  dependent. Encoder and decoder are shared and context blind, so a context
  difference has nowhere else to go.}
  \label{fig:architecture}
\end{figure}

\subsection{Model and optimisation}

Encoder and decoder are 3-layer MLPs, 192 hidden units, layer normalisation,
GELU. The context encoder is a 2-layer $\tanh$ MLP with 64 hidden units and a
small-scale output head, followed by anchoring at the reference context. The
latent dimension is 6, except in the motor simulation where it is 8 to
accommodate the condition-independent component, and $K=6$ generators are used
throughout. AdamW, learning rate $2\times10^{-3}$, weight decay $10^{-5}$, batch
size 256, gradient clipping at norm 1, cosine annealing over 200--250 epochs.
The ablations use a shorter schedule of 120 epochs on 1500 samples, applied
identically to every configuration including the unablated model, since what is
being compared there is the effect of removing an ingredient rather than the
absolute level.
Loss weights are $\lambda_{\mathrm{t}}=1$, $\lambda_{\mathrm{i}}=1$,
$\lambda_{\mathrm{g}}=0.2$, $\lambda_{\mathrm{c}}=0.05$,
$\lambda_{\mathrm{h}}=0.5$. The gauge-dependent terms are ramped in linearly over
the first 15\% of training, so the encoder is not asked to satisfy a
transformation before it has a latent to transform. Since $\theta$ depends only
on the context, the matrix exponential is evaluated once per context per step
rather than once per sample; transport and group terms are then batched over
context pairs.

For the flow gauge, $V_k(z)=G_kz+v_k+\eta\,m_k(z)$ with $\eta=0.5$ and $m_k$ a
2-layer $\tanh$ MLP whose output layer is initialised to zero, so training starts
exactly at the linear gauge. Flows use six RK4 steps. The vector-field bracket
needs $K^2$ Jacobian-vector products, which dominates the cost, so structure
constants are cached and refreshed every 25 steps and enter the objective as
derived constants, in the spirit of a target network.

\subsection{Simulations}

\textit{Place cells.} 120 cells with Gaussian fields on a sunflower tiling of a
unit disc, width $0.22$ with log-normal jitter, log-normal peak rates. Positions
come from a momentum random walk with specular reflection at the wall, which
gives the correlated, non-uniform occupancy of real foraging. Contexts act on
position through $I$, $R(60^\circ)$, $R(-40^\circ)$ with a translation, an
anisotropic stretch-and-shear $\mathrm{diag}(1.25,0.80)$ with shear $0.22$, and a
radially graded twist $u \mapsto R(0.9\,(\|u\|/R)^2)\,u$, which is a
diffeomorphism of the disc but not affine.

\textit{Grid cells.} 100 cells in one module, spacing $0.42$, orientation
$0.21$\,rad. Rates follow the standard three-cosine construction through an
exponential non-linearity. Since $k_3 = k_2-k_1$ for a hexagonal lattice, the
module's activity depends on position only through the phase pair, which is what
makes the manifold a torus; we report the homology of a single module, matching
\citet{gardner2022toroidal}. Contexts act on the phase: translations by fixed,
well-spread offsets; a $60^\circ$ lattice rotation, whose matrix
$K R_{60}K^{-1}$ is exactly integer and hence a genuine torus automorphism; and a
non-integer rescaling by $1.37$ as the control.

\textit{Motor cortex.} A 128-unit continuous-time $\tanh$ rate RNN,
$\tau\dot x = -x + W\tanh(x) + Bu(t) + b + \sqrt{2\tau}\sigma\xi(t)$, with
$\tau=50$\,ms and $\mathrm{d}t=10$\,ms, trained by BPTT for 1500 steps to emit
hand velocity toward eight targets at two extents after a go cue, following a
300\,ms preparatory period. Reaches are slightly curved, with lateral velocity
equal to the derivative of the speed profile. The curvature is there for a
reason. With perfectly straight reaches the canonical reach-plane coordinate is
confined to a line, the second row of any readout from latent to reach plane is unconstrained
by the reference context, and no method can be scored on recovering a rotation of
that plane. Recorded activity is a random dense mixture of unit rates, as with an
electrode array; twelve trials per condition are simulated with independent
intrinsic noise. Splits are taken at the level of whole trials so no trajectory
straddles the boundary.

\subsection{Baselines}

Each baseline supplies an observed-frame latent per sample per context and an
estimated linear map from the canonical frame to each context's frame; the
canonical latent is $z_c = T_c^{-1}w_c$ and every metric applies unchanged. PCA,
UMAP, the autoencoder and the VAE embed pooled activity and have their
transformations fitted post hoc by least squares on the training split.
Procrustes runs PCA per context and aligns to the reference by orthogonal
Procrustes \citep{schonemann1966generalized}. CCA aligns each context to the
reference using the paired correspondence explicitly. Manifold alignment builds a
joint graph with within-context $k$-nearest-neighbour edges and across-context
correspondence edges, embeds with the graph Laplacian \citep{ham2005semisupervised},
and extends out of sample by random-feature ridge regression.

For the headline transport score, every method including ours is read out through
the same non-linear map: a random-Fourier-feature ridge regression
\citep{rahimi2007random} from observed-frame latent to activity, fitted on the
training split. This separates the contribution of the latent and the
transformation from decoder capacity. Where a method has a native decoder we
report its own numbers too, and they are similar.

\subsection{Metric details}

Persistence diagrams come from \texttt{ripser} \citep{tralie2018ripser} on
subsamples of up to 500 points, with filtration values normalised by the mean
pairwise distance so diagrams are comparable across embeddings of different
scale. The essential class is given the diameter of the cloud as its death time;
capping it at the largest finite death instead would make it indistinguishable
from the longest noise bar and destroy the gap that Betti estimation relies on.
Before comparing diagrams we drop bars shorter than 1\% of the longest and keep
at most the forty longest, which removes the near-diagonal noise band that would
otherwise dominate an optimal matching. Betti numbers come from the largest
multiplicative gap in the sorted lifetimes, required to exceed 1.5. That
threshold is the one genuinely tuned constant in the analysis: at the sample sizes
used here a disc produces spurious $H_1$ gaps of about 1.1 and a torus produces a
true two-bar gap of about 1.9, so 1.5 separates them, but the separation for
$H_1$ on a torus is the marginal case for any gap-based estimator and we report
the achieved gap alongside the estimate.

The post-hoc algebra fitted to a baseline's transformations uses the principal
real matrix logarithm. We record the discarded imaginary part, which is large
when a transformation lies outside the principal branch, for instance a rotation
by more than $\pi$; in that case the extracted algebra element is not unique. The
GRE chart is ridge regularised with a relative penalty, which matters when the
latent dimension exceeds that of the ground-truth coordinate. An unpenalised fit
places weights of order $1/\sigma$ on low-variance latent directions, so the
chart is accurate on the reference context but extrapolates wildly onto the
transported latents it is then applied to, penalising a method for having a rich
latent rather than for getting the transformation wrong.

\clearpage
\section{Additional results}
\label{app:extra}

\begin{figure}[t]
  \centering
  \gndfig{fig6_ablations.pdf}{\textwidth}
  \caption{\textbf{Ablations.} (a) Change in transport $R^2$ when one ingredient
  is removed. (b) The same ablations on two geometric metrics, which respond
  differently. (c) Latent dimension sweep.}
  \label{fig:ablations}
\end{figure}

\begin{figure}[t]
  \centering
  \gndfig{fig7_robustness.pdf}{\textwidth}
  \caption{\textbf{Robustness and scaling.} (a,b) Observation noise. (c) Number of
  recorded cells. (d) Number of behavioural samples.}
  \label{fig:robustness}
\end{figure}

\gndinput{generated/table_grid}

\gndinput{generated/table_motor}

\gndinput{generated/table_continuous}

\gndinput{generated/table_ablations}

\gndinput{generated/table_hippocampus_per_context}

\gndinput{generated/table_grid_families}

\gndinput{generated/table_noise}
